%!TeX root=../thesis.tex

% ============================================================
%  Appendix B — Registered Hypotheses Data (Anonymised)
%  Master's Thesis: Evaluating and Refining Hypo-Stage
%  José Gonçalves Lima Neto — IME-USP, 2026
% ============================================================

\chapter{Registered Hypotheses Data}
\label{app:hypotheses-data}

This appendix presents the complete set of thirteen architectural hypotheses
registered by the two participating teams in Hypo-Stage during the observation
period. All entity references, service names, and internal artefact links have
been anonymised in accordance with the ethical agreement described in
Section~\ref{sec:ethical-considerations}. The anonymisation legend is provided
in Table~\ref{tab:entity-legend}. Technical planning entries are reproduced in
full in Section~\ref{sec:app-technical-plans}.

% ─────────────────────────────────────────────────────────────
\section{Anonymisation Legend}
\label{sec:app-legend}
% ─────────────────────────────────────────────────────────────

Hypotheses in Hypo-Stage are linked to components in the organization's Backstage
catalog through \texttt{entityRef} strings (human-readable labels that identify a
service, worker, platform capability, or template entry). To protect confidentiality
under the ethical agreement summarised in Section~\ref{sec:ethical-considerations},
this appendix does not reproduce real catalog identifiers. Instead, each hypothesis
table in Section~\ref{sec:app-hypothesis-registers} lists \emph{anonymised} entity
references in its \textbf{Entities} row; Table~\ref{tab:entity-legend} explains what
each anonymised string denotes at a functional level, without recovering the original
catalog paths or service names.

The table is intentionally typeset here as \emph{non-floating} material so it remains
with this section even though it occupies most of a page: the goal is to keep the
legend immediately available when reading the hypothesis registers, rather than
deferring it past the large tables in Section~\ref{sec:app-hypothesis-registers}.

\noindent\begin{minipage}{\linewidth}
\centering\small
\setlength{\tabcolsep}{5pt}
\begin{tabular}{p{5.2cm}p{7.8cm}}
\toprule
\textbf{Anonymised entity reference} & \textbf{Description} \\
\midrule
\texttt{catalog: Kotlin + JOOQ + Spring service template}
  & Internal catalog entry for the JVM service scaffold using Kotlin, JOOQ, and
    Spring — intended to replace a legacy functional-language stack in the
    regulated financial domain. \\[4pt]
\texttt{internal-framework: JVM platform stack with Kafka integration}
  & The organization's shared application framework for JVM services, covering
    Kafka exposure conventions, request logging, and platform integrations. \\[4pt]
\texttt{service: secure blob storage — API}
  & HTTP API for storing and governing binary documents with security and
    lifecycle rules. \\[4pt]
\texttt{service: secure blob storage — worker}
  & Background worker for the secure blob storage domain (async processing and
    maintenance). \\[4pt]
\texttt{service: confidential records \& credit-bureau integration — API}
  & API for identity, credit-bureau, and confidential document flows, currently
    coupled to a single business vertical. \\[4pt]
\texttt{service: confidential records \& credit-bureau integration — worker}
  & Worker for the same confidential records and bureau integration context. \\[4pt]
\texttt{worker: high-volume billing aggregation}
  & Long-running worker that aggregates large billing datasets. \\[4pt]
\texttt{platform: shared streaming infrastructure}
  & Shared Kafka streaming platform and consumption standards, modelled as a
    catalog component. \\
\bottomrule
\end{tabular}
\captionof{table}{Anonymised Backstage entity references and their descriptions.}
\label{tab:entity-legend}
\end{minipage}
\par\bigskip

% ─────────────────────────────────────────────────────────────
\section{Hypothesis Registers}
\label{sec:app-hypothesis-registers}
% ─────────────────────────────────────────────────────────────

Each hypothesis is typeset as a single register table: fixed metadata rows, a
full-width statement row, then author notes and technical-plan references (see
Section~\ref{sec:app-technical-plans}).

% #1 ID … #8 Uncertainty/Impact; #9 statement. Then call \hypoboxtail{notes}{plans}.
\newcommand{\hypoboxnine}[9]{%
  \par\addvspace{0.55\baselineskip}%
  \noindent\begingroup
  \small
  \setlength{\tabcolsep}{6pt}%
  \begin{tabular}{@{} >{\bfseries\raggedright\arraybackslash}p{3.4cm} >{\raggedright\arraybackslash}p{\dimexpr\textwidth-3.4cm-6\tabcolsep} @{}}
  \toprule
  \multicolumn{2}{@{}l@{}}{\large #1} \\
  \midrule
  Team & #2 \\
  Registered & #3 \\
  Lifecycle status & #4 \\
  Source type & #5 \\
  Uncertainty / Impact & #8 \\
  Quality attributes & #7 \\
  Entities & {\begingroup\ttfamily\small #6\endgroup} \\
  \addlinespace
  \multicolumn{2}{@{}>{\raggedright\arraybackslash}p{\dimexpr\textwidth-4\tabcolsep}@{}}{%
    \textbf{Statement}\par\vspace{0.3\baselineskip}%
    ``#9''\par
  } \\
  \midrule
}
\newcommand{\hypoboxtail}[2]{%
  Notes & #1 \\
  Technical plans & #2 \\
  \bottomrule
  \end{tabular}%
  \endgroup
  \par
}

% ── H01 ─────────────────────────────────────────────────────
\hypoboxnine%
  {H01}%
  {Product}%
  {2026-03-18}%
  {Open}%
  {Requirements}%
  {catalog: Kotlin + JOOQ + Spring service template}%
  {Performance \textbar{} Reliability \textbar{} Maintainability}%
  {Very High / Very High}%
  {The Kotlin + JOOQ + Spring service template is ready to scale for use in
   regulated financial product services, replacing a legacy functional-language
   stack.}%
\hypoboxtail%
  {The organization intends to deprecate the legacy stack in that domain; the
   Kotlin + JOOQ + Spring template is the leading candidate to replace it.}%
  {None registered at end of observation period.}

% ── H02 ─────────────────────────────────────────────────────
\hypoboxnine%
  {H02}%
  {Platform}%
  {2026-02-11}%
  {Open}%
  {Requirements}%
  {internal-framework: JVM platform stack with Kafka integration}%
  {Configurability \textbar{} Flexibility \textbar{} Extensibility \textbar{} Modifiability}%
  {Very Low / Medium}%
  {Kafka should be exposed to application layers for direct use via the internal
   platform framework.}%
\hypoboxtail%
  {Subject on hold; technical plans expected for the following quarter.}%
  {None registered at end of observation period.}

% ── H03 ─────────────────────────────────────────────────────
\hypoboxnine%
  {H03}%
  {Platform}%
  {2026-02-19}%
  {Open}%
  {Requirements}%
  {internal-framework: JVM platform stack with Kafka integration}%
  {Observability \textbar{} Auditability}%
  {Medium / Very High}%
  {Decouple engineering/observability logs from security/compliance auditing by
   evolving database logging toward a unified Request Logs capability on the
   internal platform framework.}%
\hypoboxtail%
  {The platform framework must let applications derive full log context including
   database logs. It is expected to supersede a legacy audit library and satisfy
   organization-wide security audit requirements.}%
  {TP-03-A, TP-03-B, TP-03-C (see Section~\ref{sec:app-technical-plans}).}

% ── H04 ─────────────────────────────────────────────────────
\hypoboxnine%
  {H04}%
  {Product}%
  {2026-03-18}%
  {Open}%
  {Requirements}%
  {service: secure blob storage — API \textbar{} service: secure blob storage — worker}%
  {Auditability \textbar{} Usability}%
  {Medium / Very Low}%
  {The secure blob storage service must retain a reference to the data subject in
   its domain data so all sensitive blobs can be deleted when the user exercises
   their right to erasure under applicable privacy regulation.}%
\hypoboxtail%
  {Concern raised in an internal RFC about fully erasing user sensitive data,
   including binary objects.}%
  {None registered at end of observation period.}

% ── H05 ─────────────────────────────────────────────────────
\hypoboxnine%
  {H05}%
  {Product}%
  {2026-03-18}%
  {In Review}%
  {Solution}%
  {service: confidential records \& credit-bureau integration — API \textbar{} \ldots{} worker}%
  {Interoperability \textbar{} Extensibility \textbar{} Reusability}%
  {High / Very High}%
  {Credit-bureau integrations can be decoupled from vertical-specific domain
   entities into context-agnostic middleware that accepts only each provider's
   minimum required data (e.g., national ID), reusable across domains without
   duplicated cost.}%
\hypoboxtail%
  {Discovery focused on: (1)~dependency matrix --- current vs.\ real dependency
   per integration; (2)~agnostic interface --- contract without domain concepts,
   only minimal identifiers and bureau outputs; (3)~migration plan --- which flows
   can isolate immediately vs.\ requiring deeper business-logic refactors.}%
  {TP-05-A (see Section~\ref{sec:app-technical-plans}).}

% ── H06 ─────────────────────────────────────────────────────
\hypoboxnine%
  {H06}%
  {Product}%
  {2026-02-24}%
  {Open}%
  {Solution}%
  {worker: high-volume billing aggregation}%
  {Performance \textbar{} Reliability \textbar{} Scalability}%
  {Very Low / Very High}%
  {Chunked keyset pagination with isolated repository-level transactions keeps
   memory strictly bounded by chunk size and prevents out-of-memory failures when
   aggregating millions of billing records.}%
\hypoboxtail%
  {Validated with a capacity/volume test at 5$\times$ production scale (45M rows
   total, 1.5M rows/month). Profiling showed JPA streams can cause first-level
   cache growth; chunking scopes transactions ($\approx$10k rows) so the garbage
   collector can reclaim memory (sawtooth heap pattern). Composite index on
   \texttt{(billing\_year\_month, status, id)}.}%
  {TP-06-A (see Section~\ref{sec:app-technical-plans}).}

% ── H07 ─────────────────────────────────────────────────────
\hypoboxnine%
  {H07}%
  {Product}%
  {2026-04-07}%
  {Validated}%
  {Requirements}%
  {service: secure blob storage — API \textbar{} service: secure blob storage — worker}%
  {Scalability \textbar{} Usability \textbar{} Configurability \textbar{} Flexibility}%
  {Very Low / Very High}%
  {Secure blob storage must support transferring blob ownership across service
   domains so blobs can move from one owning service to another.}%
\hypoboxtail%
  {Validated through a technical discovery spike that defined the transfer
   capability and produced a prioritised implementation backlog.}%
  {TP-07-A (see Section~\ref{sec:app-technical-plans}).}

% ── H08 ─────────────────────────────────────────────────────
\hypoboxnine%
  {H08}%
  {Product}%
  {2026-03-18}%
  {Validated}%
  {Requirements}%
  {service: secure blob storage — API}%
  {Security \textbar{} Usability \textbar{} Interoperability \textbar{} Extensibility}%
  {Very Low / Very High}%
  {Secure blob storage must support sharing a blob across multiple clients so the
   same object can be read from different domain services.}%
\hypoboxtail%
  {Example: a customer-facing channel receives documents later routed to support
   tooling; support needs read access in chat history without breaking vault
   invariants.}%
  {TP-08-A (see Section~\ref{sec:app-technical-plans}).}

% ── H09 ─────────────────────────────────────────────────────
\hypoboxnine%
  {H09}%
  {Platform}%
  {2026-04-12}%
  {Open}%
  {Solution}%
  {platform: shared streaming infrastructure}%
  {Auditability \textbar{} Reliability}%
  {Medium / Very High}%
  {Using the Dead-Letter Topic (DLT) only for auditing and diagnosis, and
   reprocessing from the original topic's offset, preserves consistency and
   auditability better than re-enqueueing messages from the DLT.}%
\hypoboxtail%
  {How to falsify: controlled replay from the DLT without mutating the event must
   be proven safer, simpler, and more auditable than replay by offset.}%
  {TP-09-A, TP-09-B (see Section~\ref{sec:app-technical-plans}).}

% ── H10 ─────────────────────────────────────────────────────
\hypoboxnine%
  {H10}%
  {Platform}%
  {2026-04-12}%
  {Open}%
  {Solution}%
  {platform: shared streaming infrastructure}%
  {Interoperability \textbar{} Performance}%
  {Medium / High}%
  {For short-lived transient failures, limited synchronous retry on the consumer
   thread preserves per-partition ordering better and with lower overall complexity
   than asynchronous retry on a separate queue.}%
\hypoboxtail%
  {How to falsify: a per-partition retry queue preserves order with less blocking
   time and acceptable operational cost.}%
  {TP-10-A (see Section~\ref{sec:app-technical-plans}).}

% ── H11 ─────────────────────────────────────────────────────
\hypoboxnine%
  {H11}%
  {Platform}%
  {2026-04-12}%
  {Open}%
  {Solution}%
  {platform: shared streaming infrastructure}%
  {Reliability \textbar{} Modifiability \textbar{} Interoperability}%
  {Medium / High}%
  {Keeping failure state inside the Kafka message flow increases coordination
   complexity across threads and complicates failure handling, compared to
   handling failures in the execution layer and keeping consumption non-blocking.}%
\hypoboxtail%
  {How to falsify: modelling failure state in Kafka shows lower operational
   complexity and equal or better recovery than handling failures in the execution
   layer.}%
  {TP-11-A (see Section~\ref{sec:app-technical-plans}).}

% ── H12 ─────────────────────────────────────────────────────
\hypoboxnine%
  {H12}%
  {Platform}%
  {2026-04-12}%
  {Open}%
  {Solution}%
  {platform: shared streaming infrastructure}%
  {Performance \textbar{} Reliability}%
  {Medium / High}%
  {Classifying errors as recoverable vs.\ non-recoverable before applying retry
   reduces consumer lag and computational cost without worsening the successful
   processing rate.}%
\hypoboxtail%
  {How to falsify: indiscriminate retries improve final success rate without
   meaningful latency/capacity cost.}%
  {TP-12-A, TP-12-B (see Section~\ref{sec:app-technical-plans}).}

% ── H13 ─────────────────────────────────────────────────────
\hypoboxnine%
  {H13}%
  {Platform}%
  {2026-04-12}%
  {Open}%
  {Solution}%
  {platform: shared streaming infrastructure}%
  {Auditability}%
  {Very Low / Very High}%
  {Allowing fixes by mutating the payload before re-enqueueing increases
   consistency violations and makes traceability and auditing harder.}%
\hypoboxtail%
  {How to falsify: a formal correction mechanism fully preserves audit trail,
   causality, and reproducibility of the event.}%
  {TP-13-A (see Section~\ref{sec:app-technical-plans}).}

\bigskip
\section{Technical Plans Catalog}
\label{sec:app-technical-plans}
% ─────────────────────────────────────────────────────────────

The following entries reproduce all technical plans registered in Hypo-Stage
during the observation period. Each entry is identified by a plan code
(e.g., TP-03-A), the parent hypothesis, the action type, the target date, a
description of the planned action, and the expected outcome. Internal artefact
links (tickets, pull requests, design documents) have been redacted.

\vspace{6pt}

%% ─── H03 plans ────────────────────────────────────────────
\paragraph{TP-03-A \textnormal{(H03 — Decouple eng.\ logs from security/compliance audit)}}
\textit{Action type:} Architectural Spike \quad \textit{Target date:} 2026-03-31 \\
\textit{Description:} Discover and implement default mapping of database
transaction IDs into the Request Log structure within the internal framework
so every request and application log entry carries database transaction IDs for
correlation and audit trail without per-application integration. \\
\textit{Expected outcome:} Request logs consistently include database transaction
ID fields; applications using the framework obtain audit-ready context without
custom wiring.

\paragraph{TP-03-B \textnormal{(H03)}}
\textit{Action type:} Modularisation / Flexibility \quad \textit{Target date:} 2026-03-31 \\
\textit{Description:} Extend framework logging to worker modules so critical
database transactions in workers log with the same context (transaction IDs,
correlation) as request paths. \\
\textit{Expected outcome:} Workers emit logs with database transaction IDs
aligned with the Request Log infrastructure; audit coverage for
asynchronous/background work.

\paragraph{TP-03-C \textnormal{(H03)}}
\textit{Action type:} Modularisation / Flexibility \quad \textit{Target date:} 2026-03-31 \\
\textit{Description:} Stream database transaction properties to the same log
buckets and pipeline as Request and Application logs for a single
audit/observability backbone. \\
\textit{Expected outcome:} Database transaction properties follow the same
ingestion and retention path as Request logs; no separate audit-only pipeline.

%% ─── H05 plan ─────────────────────────────────────────────
\paragraph{TP-05-A \textnormal{(H05 — Credit-bureau decoupling seam)}}
\textit{Action type:} Discovery \quad \textit{Target date:} ongoing \\
\textit{Description:} Structured discovery covering: (1)~dependency matrix
(current vs.\ real dependencies per integration); (2)~agnostic interface
proposal (contract free of domain concepts, accepting only minimal provider
identifiers and bureau outputs); (3)~migration sequencing (which flows can
isolate immediately vs.\ requiring deeper business-logic refactors). \\
\textit{Expected outcome:} A prioritised integration decoupling plan and a
draft interface contract.

%% ─── H06 plan ─────────────────────────────────────────────
\paragraph{TP-06-A \textnormal{(H06 — Billing aggregation chunked keyset pagination)}}
\textit{Action type:} Implementation (Other) \quad \textit{Target date:} 2026-03-13 \\
\textit{Description:} Implement the billing aggregation worker with chunked
keyset pagination using a \texttt{LIMIT} loop with
\texttt{id > :lastProcessedId}; avoid broad \texttt{@Transactional} so the
Hibernate first-level cache clears per chunk; add composite index on
\texttt{(billing\_year\_month, status, id)}. \\
\textit{Expected outcome:} Process 300k+ rows/month safely without
out-of-memory failures; stable sawtooth heap profile; healthy database
CPU/I/O via index scans.

%% ─── H07 plan ─────────────────────────────────────────────
\paragraph{TP-07-A \textnormal{(H07 — Blob ownership transfer)}}
\textit{Action type:} Architectural Spike \quad \textit{Target date:} 2026-04-30 \\
\textit{Description:} Technical discovery to define the transfer capability on
the blob storage API. \\
\textit{Expected outcome:} A backlog for implementing the ownership transfer
feature.

%% ─── H08 plan ─────────────────────────────────────────────
\paragraph{TP-08-A \textnormal{(H08 — Blob cross-client sharing)}}
\textit{Action type:} Guideline Implementation \quad \textit{Target date:} 2026-04-06 \\
\textit{Description:} Alignment session on cross-domain sharing and transfers
so that long-lived access (e.g., messaging history) does not require definitive
ownership by every consumer. \\
\textit{Expected outcome:} A governance decision enabling prioritisation of
sharing and transfer implementation work.

%% ─── H09 plans ────────────────────────────────────────────
\paragraph{TP-09-A \textnormal{(H09 — DLT for audit vs.\ offset replay)}}
\textit{Action type:} Guideline Implementation \quad \textit{Target date:} 2026-06-30 \\
\textit{Description:} Produce a runbook for reprocessing by offset with minimal
tooling for controlled replay. \\
\textit{Expected outcome:} Documentation of how to reprocess by offset with
minimal tooling.

\paragraph{TP-09-B \textnormal{(H09)}}
\textit{Action type:} Experiment \quad \textit{Target date:} 2026-06-30 \\
\textit{Description:} Operational exercise with a simulated incident. \\
\textit{Expected outcome:} Measure MTTR, duplicate rate, and divergence between
the original event and the final state.

%% ─── H10 plan ─────────────────────────────────────────────
\paragraph{TP-10-A \textnormal{(H10 — Synchronous retry for transient failures)}}
\textit{Action type:} Tracer Bullet \quad \textit{Target date:} 2026-06-30 \\
\textit{Description:} Tracer bullet on a critical consumer path; inject
transient failures with windows of 1s, 5s, and 30s. \\
\textit{Expected outcome:} Measure out-of-order delivery, consumer lag,
throughput, and retry success rate.

%% ─── H11 plan ─────────────────────────────────────────────
\paragraph{TP-11-A \textnormal{(H11 — Failure state inside vs.\ outside Kafka flow)}}
\textit{Action type:} Architectural Spike \quad \textit{Target date:} 2026-06-30 \\
\textit{Description:} Spike comparing both failure-state approaches on a real
consumer. \\
\textit{Expected outcome:} Measure consumer lag, blocked-thread time,
states/transitions, and operational recovery steps for each approach.

%% ─── H12 plans ────────────────────────────────────────────
\paragraph{TP-12-A \textnormal{(H12 — Error classification before retry)}}
\textit{Action type:} Guideline Implementation \quad \textit{Target date:} 2026-06-30 \\
\textit{Description:} Define an error taxonomy and mandatory classification step
before retry is applied. \\
\textit{Expected outcome:} Practical directives for classifying errors before
retry, consumable by all teams using the shared streaming platform.

\paragraph{TP-12-B \textnormal{(H12)}}
\textit{Action type:} Software Analytics (Architectural Trigger) \quad \textit{Target date:} 2026-06-30 \\
\textit{Description:} Instrument analytics by error class on a real consumer:
success rate, average attempts, extra time per message, and discards to DLT. \\
\textit{Expected outcome:} A metrics dataset sufficient to falsify or support
the hypothesis.

%% ─── H13 plan ─────────────────────────────────────────────
\paragraph{TP-13-A \textnormal{(H13 — Payload immutability before re-enqueueing)}}
\textit{Action type:} Experiment \quad \textit{Target date:} 2026-06-30 \\
\textit{Description:} Enforce an immutable-payload guideline; add automated
validation in replay tooling; require mandatory architectural review for
exceptions. \\
\textit{Expected outcome:} Experiment results that validate or refute the
hypothesis and establish a baseline for payload governance in the shared
streaming platform.
